Structured prediction models can leverage correlations that allow them to make correct predictions even with very little underlying evidence.
%between output variables, which can potentially result in bias and, as we showed, amplified bias.
Yet such models risk potentially leveraging social bias in their training data.
In this paper, we presented a general framework for visualizing and quantifying biases in such models and proposed \alg to calibrate their predictions under two different settings.
Taking gender bias as an example, our analysis demonstrates that conditional random fields can amplify social bias from data while our approach \alg{} can help to reduce the bias.

Our work is the first to demonstrate structured prediction models amplify bias and the first to propose methods for reducing this effect but significant avenues for future work remain.
While \alg{} can be applied to any structured predictor, it is unclear whether different predictors amplify bias more or less.
Furthermore, we presented only one method for measuring bias. 
More extensive analysis could explore the interaction among predictor, bias measurement, and bias de-amplification method. 
Future work also includes applying bias reducing methods in other structured domains, such as pronoun reference  resolution~\cite{mitkov2014anaphora}.


\paragraph{Acknowledgement}
This work was supported in part by National Science Foundation Grant IIS-1657193 and two NVIDIA Hardware Grants.